\section{Methodology}
\label{sec:methodology}

\subsection{Hypothesis}

\textbf{H1:} A state database based on Sparse Merkle Tree with
Poseidon hash implemented in Go can support 10,000+ accounts with
state root computation $< 50$~ms, compatible with witness generation
for the ZK prover.

\textbf{H0 (null):} A Go SMT with Poseidon cannot achieve state root
computation $< 50$~ms at 10,000+ accounts, or the implementation
introduces incompatibility with the BN254 field arithmetic required
by the ZK prover, or memory usage exceeds 1~GB at 10,000 accounts.

\subsection{Predictions}

Prior to experimentation, we registered the following quantitative
predictions based on published benchmarks and the RU-V1 baseline:

\begin{enumerate}
  \item Go Poseidon2 hash will be 50--200$\times$ faster than
    JavaScript BigInt due to native field arithmetic.
  \item Insert latency will be $< 100$~$\mu$s per account
    (vs.\ 1.8~ms in TypeScript RU-V1).
  \item Proof generation will be $O(d)$ traversal, $< 10$~$\mu$s
    per proof.
  \item State root computation for 10,000 accounts via batch insert
    will complete in $< 50$~ms.
  \item Memory usage will be $< 500$~MB at 10,000 accounts with
    in-memory storage.
\end{enumerate}

\subsection{Experimental Design}

The experiment follows a controlled benchmarking methodology with
two implementation phases:

\textbf{Phase 1: Baseline implementation.} SMT using Go's
\texttt{math/big.Int} for field elements, wrapping gnark-crypto's
Poseidon2 permutation. This establishes correctness and provides
initial performance measurements.

\textbf{Phase 2: Optimized implementation.} SMT using gnark-crypto's
native \texttt{fr.Element} type (Montgomery form, assembly-optimized),
eliminating \texttt{big.Int} allocation overhead. This represents the
production-viable implementation.

Both phases use identical tree parameters (depth 32, BN254 scalar
field) and benchmark workloads (100, 1,000, and 10,000 entries)
to enable direct comparison.

\subsection{Variables}

\textbf{Independent variables:}
\begin{itemize}
  \item Account count: 100, 1,000, 10,000
  \item Tree depth: 32 (matching RU-V1 baseline)
  \item Batch size: 10, 50, 100, 250, 500 transactions
  \item Implementation variant: \texttt{big.Int} vs.\ \texttt{fr.Element}
\end{itemize}

\textbf{Dependent variables:}
\begin{itemize}
  \item Hash latency ($\mu$s per Poseidon2 2-to-1 compression)
  \item Insert latency ($\mu$s per SMT insert operation)
  \item Proof generation latency ($\mu$s per Merkle proof)
  \item Proof verification latency ($\mu$s per proof verification)
  \item Batch update latency (ms per batch of $n$ updates)
  \item Memory usage (MB at $n$ entries)
\end{itemize}

\textbf{Controls:}
\begin{itemize}
  \item Field: BN254 scalar field ($p = 2^{254} + \epsilon$)
  \item Hash: Poseidon2 with $r_f = 6$, $r_p = 50$
  \item Benchmark repetitions: 50 per configuration
  \item Warmup iterations: 10 (excluded from measurements)
  \item Key generation: deterministic sequential keys for
    reproducibility
\end{itemize}

\subsection{Metrics and Evaluation Criteria}

All metrics are standard in ZK systems literature:

\begin{itemize}
  \item \textbf{Latency metrics} reported as mean, standard
    deviation, and percentiles (P50, P95, P99) over 50 repetitions.
  \item \textbf{Throughput} derived as inverse of per-operation
    latency.
  \item \textbf{Memory} measured via Go's \texttt{runtime.MemStats}
    after garbage collection.
  \item \textbf{Speedup} computed as ratio of TypeScript RU-V1
    baseline to Go optimized implementation.
\end{itemize}

The hypothesis is \textbf{confirmed} if batch update latency for
a realistic block size ($\leq 250$ transactions) on a 10,000-entry
tree remains below 50~ms, with memory under 1~GB.
